% §403 `scan_left_brace` tests the *command*, not the character:
%
%     procedure scan_left_brace; {reads a mandatory left brace}
%     begin @<Get the next non-blank non-relax non-call token@>;
%     if cur_cmd<>left_brace then
%       begin print_err("Missing { inserted");
%       help4("A left brace was mandatory here, so I've put one in.")…
%
% `cur_cmd = left_brace` is true of an explicit `{` and equally of a control
% sequence `\let` to one — plain's `\bgroup`. Everything that reads a mandatory
% brace goes through here: `\hbox`, `\vbox`, `\vtop`, `\vcenter`, the three
% parts of `\discretionary`, `\insert`, `\vadjust`, `\output`, and §1151's
% `scan_math` for a subformula.
%
% mathtools builds every paired delimiter on the implicit form:
%
%     \DeclarePairedDelimiter{\norm}{\lVert}{\rVert}
%
% expands `\norm{Z}` through `\mathopen{}\mathclose\bgroup … \aftergroup\egroup`,
% so the brace that opens the subformula is `\bgroup` and never a `{`. Reading
% the token's catcode instead of its meaning refuses it:
%
%     ! Missing { inserted.
%     <recently read> \delimiter
%     l.321  \norm{Z}
%                    _{\psi_2}:=\inf\left\{r>0:
%
% 2608.15840v1 uses `\norm` throughout and stopped at the hundred-error limit
% after four of its thirty-six pages.
\catcode`\{=1 \catcode`\}=2 \catcode`\$=3 \catcode`\^=7 \catcode`\_=8
\font\r=cmr10 \r \hsize=200pt \parindent=0pt
\font\mi=cmmi10 \font\sy=cmsy10 \font\ex=cmex10
\textfont0=\r \scriptfont0=\r \scriptscriptfont0=\r
\textfont1=\mi \scriptfont1=\mi \scriptscriptfont1=\mi
\textfont2=\sy \scriptfont2=\sy \scriptscriptfont2=\sy
\textfont3=\ex \scriptfont3=\ex \scriptscriptfont3=\ex
\let\bgroup={ \let\egroup=}

% A box whose brace is implicit — the plainest case of §403.
\setbox0=\hbox\bgroup A\egroup
\message{<hbox=\the\wd0>}
\setbox0=\vbox\bgroup \hbox{B}\egroup
\message{<vbox=\the\ht0>}
\setbox0=\vtop\bgroup \hbox{C}\egroup
\message{<vtop=\the\ht0>}

% §1151's `scan_math`: a subformula opened by an implicit brace. This is the
% one mathtools depends on, reached through `\mathclose\bgroup`.
\setbox0=\hbox{$x\mathclose\bgroup y\egroup$}
\message{<mathclose=\the\wd0>}
\setbox0=\hbox{$\mathopen\bgroup z\egroup$}
\message{<mathopen=\the\wd0>}

% A script field, which also reads a mandatory brace.
\setbox0=\hbox{$a^\bgroup b\egroup$}
\message{<script=\the\wd0>}

% `\vcenter`, and a `\vbox to` where the dimension is scanned before the brace.
\setbox0=\hbox{$\vcenter\bgroup\hbox{D}\egroup$}
\message{<vcenter=\the\wd0>}
\setbox0=\vbox to 20pt\bgroup \hbox{E}\vfil\egroup
\message{<vboxto=\the\ht0>}

% `\discretionary`'s three parts.
\setbox0=\hbox{a\discretionary\bgroup-\egroup\bgroup\egroup\bgroup\egroup b}
\message{<disc=\the\wd0>}

% An explicit brace still works, and so does a brace reached through one more
% level of `\let` — the meaning is what counts, however far it is from a `{`.
\let\obrace=\bgroup \let\cbrace=\egroup
\setbox0=\hbox\obrace F\cbrace
\message{<relet=\the\wd0>}

% A *space* before the implicit brace: §403 skips blanks first, so this is the
% same case and not a different one.
\setbox0=\hbox \bgroup G\egroup
\message{<spaced=\the\wd0>}

% The control: a token that is not a brace at all still has to be refused, or
% the fix would accept anything. `\relax` is skipped by §403's "non-relax"
% clause, so the `H` is what gets refused — and TeX inserts the brace and
% carries on, which is what makes this comparable.
\message{<end>}
\end
